\scanpage{8}
Les deux premières assertions résultent
du fait suivant : \uline{Soit} $H$ \uline{un sous-schéma
fermé de} $G$, \inserted{\uline{propre et}} \uline{plat sur} $S$. \uline{Alors il existe un}
\deleted{\uline{plus} \uline{petit}} \uline{sous-schéma fermé} $\deleted{S'}T$ \uline{de} $S$
tel que pour tt morphisme $\deleted{\unsure{T}}S'\mathrel{-}S$,
l'image inverse $H'$ est un ss-groupe de $G'$
\unsure{sss|ssi} $S'\mathrel{-}S$ est majoré par $T$. Or,
il faut exprimer, \deleted{si \unsure{$H$ et|$H$ est}} $1^\circ$) Les images
inverses de $H$ et $\check H$ sont les ${=}_{s}$
$2^\circ$) \deleted{\unsure{Les|l'} images}\inserted{L'image} inverse de $H\times H$ est
majorée par celle de $P$, où
$P=p^{-1}(H)$, $p:G\times G\mathrel{-}G$ est la
multiplication. Les deux \deleted{\unsure{que}} conditions
en question, il est connu, s'expriment
par la \deleted{\unsure{s}} domination de $S'\mathrel{-}S$ par
des ss-schémas fermés $T_1$, $T_2$ de $S$,
on prendra donc $T=T_1\unsure{\wedge|\cap} T_2$.

Pour la dernière assertion, le critère
habituel \unsure{des|de} propreté \unsure{nous} ramène au
Cor. 4 \unsure{,} de prop. 1. O.K.
% archive sheet number: 7

